\section{Setting, protocol, and evidence}
\label{sec:methods}

\subsection{Task and problem source}

The full user objective is reproduced in Appendix~\ref{app:protocol}.
It specifies an overall time limit and a research objective, rather
than a list of preselected questions or a target number of outputs.
The agent's initial plan explicitly prioritized precise statements,
literature checks, finite counterexamples, elementary reductions, and
independent implementations of computational checks. It also required
partial results and unsuccessful searches to be recorded.

The input file \artifact{docs/21tkt.pdf} has SHA-256 digest
\begin{center}
\small\path{2fcce9b98a4df10267fe120229217bfe556c70510704e311540da0cef438f911}.
\end{center}
The agent extracted a main-body index with 1,308 entries, including
all 150 entries of Issue 21. This is an index of problem statements,
not a verified count of open problems. Editorial stars, solved
subparts, duplicates, and literature published after a problem was
posed all affect the open status of individual assertions.
The source PDF, rather than its text extraction, was treated as
authoritative for mathematical notation and scope.

\subsection{Agent and computational environment}

The experiment used \texttt{gpt-6-astra} with reasoning effort
\texttt{xhigh}. These two settings were supplied by the user during
manuscript preparation on 14 September 2026; the supplied session
export does not itself retain them. The export identifies the
application as Codex 0.153.4. Official documentation identifies
\texttt{xhigh} as a supported Astra reasoning setting~\cite{astra},
but documentation of a product's capabilities is not evidence that
every such capability was exercised in this particular run.

The agent worked in a persistent Linux workspace with shell access,
file editing, Git, internet source retrieval, and tools for inspecting
PDF pages. GAP 4.16.1 supplied the main finite-group computations.
The retained scripts also use Python and C/C++ for exact arithmetic,
word calculations, independent finite models, enumeration, and
certificate parsing. The source and report files record which
libraries or catalogues each computation depends on. Versions not
retained in the experimental records cannot be reconstructed merely
by inspecting the later manuscript environment.

The user initially supplied an inaccessible GAP path and repaired it
during setup. The visible human message is \emph{sorry, gap is now here}.
The subsequent reports establish that GAP was available with the
required packages. The initial suggestion that a replacement
installation might be needed is therefore not evidence that such an
installation actually occurred.

The working limits were generally 20 CPU cores and 100\,GB of memory.
The repository contains bounded launchers and process observations,
but not a continuous, comprehensive record of global peak resource
use. In the largest retained certificate check, six GAP workers each
had a 14\,GiB workspace ceiling, for a combined configured ceiling of
84\,GiB. A workspace ceiling is not an observed resident-memory peak.
Likewise, local CPU limits do not measure the remote compute used by
the language model.

The research plan prohibited delegation to other agents without
authorization, and the reports record no delegated agents. Parallel
GAP or C++ processes were computational workers, not separate language
model researchers. The same agent developed the arguments, wrote
checkers, and conducted the internal reviews. The phrase
\emph{independent implementation} below therefore refers to a separate
representation or algorithm, not to an independent human author or
an independently trained reviewing model.

\subsection{Persistence and local records}

The workspace served as external memory. Plans and status files recorded
the next steps; proof drafts preserved mathematical arguments; scripts
and logs preserved computations; and Git recorded successive states.
The visible export contains 53 goal-continuation directives and 52
assistant final messages. These observations are consistent with a
long task continued across many turns, rather than a single uninterrupted
response. They do not reveal the full underlying context-management
implementation.

There were 198 commits during the 48-hour window. The first, at
21:00:37 UTC on 10 September, initialized the investigation.
The 199th commit, \snapshot{}, completed the local handoff at
20:57:02 UTC on 12 September. Commit times date repository checkpoints;
they do not directly date the discovery of an idea or establish when
its full final scope was proved.

At \snapshot{}, the Git tree contains 4,147 tracked files with a total
logical size of 2,440,677,642 bytes. This includes 430 files under
\artifact{research/}, 456 under \artifact{scripts/}, and 2,755 under
\artifact{results/}, as well as references, the source Notebook, and
other records. Logical file size is not compressed repository size.
The manuscript analysis reads this fixed revision, so later writing
and review do not increase the experimental counts.

\subsection{What the session export retains}

Two files were supplied after the experiment: \artifact{session.html}
and \artifact{session.txt}. Their digests, exact aggregate measurements,
and a standard-library parser are included with the manuscript's data.
The HTML header advertises 34,102 events. Its actual retained event
elements comprise 709 assistant commentaries, 52 assistant final
messages, one human message, 61 directives, and 3,616 token records:
4,439 elements in total. The text export agrees on the visible message
counts.

No tool-call bodies or internal-analysis event bodies occur as exported
event elements. Thus the export is not a complete raw execution trace,
and it cannot establish an exact tool-call count. Nor does its single
human-message element prove that there was only one human instruction:
the initial research objective is preserved inside a directive.
Most retained events have sequence identifiers but no individual
wall-clock timestamp. Narrative ordering can use these identifiers;
precise time claims require a dated repository or process record.

\subsection{Token accounting}

Table~\ref{tab:tokens} gives the final cumulative counters displayed by
the session renderer. The parser checks that the cumulative counters
are monotone, that total tokens equal input plus output, and that the
cached-input and reasoning-output counters are contained in their
respective parent counters. A second extraction of the final HTML
usage section gives the same values.

\begin{table}[ht]
\centering
\caption{Final cumulative token counters in the supplied export.
Subcounters are indented and must not be added to their parent counters.}
\label{tab:tokens}
\begin{tabular}{lr}
\toprule
Reported counter & Tokens\\
\midrule
Input & 473,526,483\\
\quad Cached input & 453,222,656\\
\quad Cache-write input & 0\\
Output & 3,807,698\\
\quad Reasoning output & 1,845,520\\
Total (input plus output) & 477,334,181\\
\midrule
Derived input minus cached input & 20,303,827\\
\bottomrule
\end{tabular}
\end{table}

Cached input accounts for approximately 95.71\% of reported input.
The 477 million total is therefore neither a count of newly generated
text nor a count of distinct tokens read for the first time.
The distinction between cached input and input, and between reasoning
output and output, follows the usage conventions described in the
official documentation~\cite{promptcaching,reasoning}.

There are 96 identical successive cumulative snapshots. In addition,
80 \emph{Last} records contain a positive total but zero for every
itemized component. The analysis retains these records without
assigning an undocumented request type to them. It does not sum
these records to reconstruct usage. These limitations, and the
absence of a reconciliation with provider invoices, prevent an
audited monetary-cost estimate from being inferred from the export
alone.

\subsection{Outcome categories and checks}

We distinguish five questions about an artifact:
\begin{enumerate}
\item Does its mathematical statement match the question being addressed?
\item Is the argument correct under its stated assumptions and imported results?
\item Do the computations establish the particular finite claims assigned to them?
\item Is the assertion already known or an immediate consequence of prior work?
\item Has someone independent of the producing agent accepted the argument?
\end{enumerate}
These questions require different evidence. A digest match answers
none of the mathematical questions. A finite check cannot establish
an unrestricted infinite-group theorem without a proved reduction.
Failure to find a paper in a literature search does not establish
priority. A second implementation can expose an error in a first
program while sharing the same mistaken mathematical specification.

The deadline ledger uses \emph{complete candidate} for an argument
presented as answering the recorded scope and surviving the internal
checks conducted during the run. A \emph{partial result} proves a
restricted assertion while leaving the full question open.
A \emph{rediscovery} reproduces a known phenomenon or proof after
independent work, whereas a \emph{prior-result deduction} applies an
identified theorem. A bounded search without a hit remains
\emph{computational evidence} for its exact range.
These are reporting categories, not interchangeable levels on a
single scale of confidence.
